\documentclass{article}
\usepackage{amsmath, amssymb, amsthm}
\newtheorem{proposition}{Proposition}
\newtheorem{theorem}{Theorem}
\newtheorem{corollary}{Corollary}
\newtheorem{definition}{Definition}
\newtheorem{remark}{Remark}
\usepackage{graphicx}
\usepackage{hyperref}
\usepackage{geometry}
\usepackage{algorithm}
\usepackage{algpseudocode}
\usepackage{xcolor}
\usepackage{float}
\usepackage{natbib}
\usepackage{booktabs}
\geometry{margin=1in}

\newcommand{\merw}{EigenTree}
\newcommand{\merwft}{EigenTree-FT}
\newcommand{\psiback}{\psi_{\mathrm{backup}}}

\title{\merw: A Decentralized Communication Protocol\\Achieving Centralized Multi-Armed Bandit Results,\\with Fault-Tolerant Hub Failover}
\author{Salvatore Rastelli, Andrea Perin}
\date{July 2026}

\begin{document}

\maketitle

\begin{abstract}
	In decentralized cooperative learning, no agent has global knowledge, and yet centralized coordination is statistically efficient.
	We propose \merw{}, a decentralized communication protocol on undirected graphs that establishes, from purely local interactions, a de-facto coordinator among the agents.
	\merw{} uses eigenvector centrality---as computed by a gossip-based power iteration---to 1) choose the most central node as coordinator, and 2) construct a routing tree, rooted at the coordinator node, which spans the graph.
	Information flows according to this routing tree: from leaves to root during the aggregation phase, and vice versa during the broadcasting phase.
	Thus, a centralized aggregator emerges out of purely decentralized interactions, at the cost of introducing appropriate communication rounds between agents in the tree.
	The resulting protocol can be directly used in various decentralized cooperative learning setups, of which we show two---regret minimization and best arm identification---where the method recovers centralized-like results.
	We further address the hub's structural weakness as a single point of failure: we propose \merwft, an extension that tolerates permanent hub failure with zero communication overhead during normal operation, by designating the hub's most central direct child as a passive backup that mirrors the hub's state for free and promotes itself upon detecting hub silence.
\end{abstract}

\section{Introduction}

In the cooperative multi-agent multi-armed bandit (MA-MAB) setting, $N$ agents are connected by an undirected graph $G = (V, E)$. At each round, every agent pulls one arm and receives a stochastic reward; agents communicate with neighbors to improve collective performance. Two objectives are studied: \emph{regret minimization}, where the goal is to minimize cumulative group regret over a horizon $T$; and \emph{Best Arm Identification} (BAI), where the goal is to identify the best arm with high confidence using as few total arm pulls as possible.

The central difficulty in both settings is the tension between \emph{decentralization} (no agent has global knowledge) and \emph{statistical efficiency} (centralized coordination is optimal). Existing approaches resolve this either by symmetric averaging (e.g., Coop-UCB2 \citep{landgren2021distributed}, which uses doubly stochastic consensus weights to diffuse information uniformly) or by assuming all-to-all communication (e.g., Hillel et al.\ \citep{hillel2013distributed}, whose distributed successive elimination requires every agent to broadcast to all others in each round).

\merw{} takes a different route: we show that a coordinator can \emph{emerge} from the graph structure itself, without any agent knowing the global topology. The key insight is that the Maximal Entropy Random Walk (MERW) eigenvector centrality, computed via a fully decentralized gossip protocol, identifies the structurally most central node. This node, the hub, then organizes the rest of the network into a spanning routing tree via max-flooding, again using only local information.

Once the tree is in place, the hub acts as a de-facto central coordinator. Every agent routes its observations up toward the hub, which sees the globally aggregated picture and makes decisions on behalf of the network. The downlink carries only what is needed: in the regret setting, the hub's enriched estimates; in the BAI setting, the surviving arm set and a timing signal so all agents restart each elimination round simultaneously.

The contribution of \merw{} is therefore twofold. First, it provides a fully decentralized protocol for \emph{inducing centrality}: the hub and tree emerge from local MERW computations, requiring no privileged node, no global communication, and no prior knowledge of the graph structure. Second, it demonstrates that this emergent coordinator is statistically sufficient: the hub achieves centralized UCB performance for regret minimization, and the same sample complexity as Hillel's protocol for BAI, while the communication overhead over Hillel is only $O(D)$ extra time steps per elimination round, where $D$ is the depth of the MERW-induced tree.

This emergent-coordinator design has a structural weakness, however: the hub $h^\star$ is a single point of failure. If it crashes permanently, no other node holds the global statistics, no downlink is broadcast, and all non-hub agents are left waiting indefinitely. Restarting the full \merw{} pipeline to elect a new hub is infeasible during an ongoing bandit episode: it requires $\tau_{\mathrm{init}}$ power iteration rounds plus a max-flooding phase, and the new hub would start from scratch with no accumulated statistics. We therefore also propose \merwft{} (Section~\ref{sec:ft}), an extension that tolerates permanent hub failure with zero overhead during normal operation, by designating the hub's most central direct child as a passive backup.

\section{Related Work}
\label{sec:related}

\paragraph{Cooperative multi-agent bandits.}
The cooperative MA-MAB problem has been studied under two broad communication models. In the \emph{gossip/consensus} model, agents exchange information with neighbors and update local estimates via averaging; no coordinator is assumed. Coop-UCB2 \citep{landgren2021distributed} uses doubly stochastic consensus weights to diffuse observations across the graph, achieving sublinear group regret that depends on the spectral gap of the communication matrix. Martinez-Rubio et al.\ \citep{martinez2019decentralized} study a similar setting and introduce the total-pull clock as a way to index UCB confidence intervals by actual observations rather than wall-clock rounds, an idea we adopt directly. These consensus-based approaches are symmetric: every agent plays an equivalent role, and there is no notion of a coordinator.

Hillel et al.\ \citep{hillel2013distributed} study cooperative BAI under a fully connected communication model, where in each round every agent may broadcast a message to all others. Their main result is a tradeoff between communication rounds and sample complexity: with a single communication round, the best achievable speedup is $\sqrt{k}$ (and this is tight by their lower bound); with $O(\log(1/\varepsilon))$ rounds, their Algorithm 3 achieves the ideal factor-$k$ speedup. We build on Algorithm 3: \merw-BAI uses the same pull schedule and elimination rule, but replaces the all-to-all broadcast with hop-by-hop relay over the MERW spanning tree. Because every agent receives the same global information after each broadcast in Hillel's setting, all agents compute the same surviving arm set independently; in our setting, the hub is the only node that sees the full aggregate and computes the elimination -- a role that emerges from the graph structure rather than being assumed by the communication model.

\paragraph{Eigenvector centrality and MERW.}
Eigenvector centrality scores nodes by their contribution to the leading eigenvector of the adjacency matrix and is widely used in network analysis. The Maximal Entropy Random Walk \citep{burda2009localization} provides an information-theoretic grounding for this measure: $\psi_i$ is the unique centrality consistent with assigning equal probability to all paths of equal length. The localization phenomenon -- that $\psi_i^2$ concentrates sharply on a small set of nodes in heterogeneous graphs -- was characterized by Burda et al.\ \citep{burda2009localization} and motivates using $\psi$ to identify a hub rather than, say, degree or betweenness centrality.

\paragraph{Decentralized eigenvector computation.}
Distributed power iteration for computing the leading eigenvector of a matrix defined by a network has been studied in the systems literature. Jelasity et al.\ \citep{jelasity2007asynchronous} give a gossip-based normalization that avoids the global all-reduce required by standard power iteration, at the cost of $O(\log N)$ extra gossip rounds per iteration. We use their protocol directly, with both the fixed iteration count and the gossip-round count set to $O(\log N)$.

\paragraph{Fault tolerance in tree-structured coordination.}
The passive backup pattern we use for \merwft{} is well-established in distributed systems. OSPF elects a Backup Designated Router (BDR) that mirrors every topology update from the Designated Router; upon DR failure the BDR promotes itself in one hello-interval \citep{rfc2328}. RSTP pre-identifies an Alternate Port for every non-root switch and activates it in a single proposal/agreement round \citep{ieee8021w}. The formal correctness of passive replication is established by Lamport, Malkhi and Zhou \citep{lamport2009vertical}, who show via reduction to Paxos that a backup mirroring the primary's state is at most one synchronization round stale and that failover is safe. \merwft{} applies this pattern to the aggregation-tree setting of \merw. The backup-selection procedure is also directly analogous to the Diffusing Update Algorithm (DUAL) used in EIGRP \citep{garcia1993eigrp} (Section~\ref{sec:eigrp}), and non-hub re-attachment builds on the swap-edge framework of Nardelli et al.\ \citep{nardelli2003} and its distributed variant by Datta et al.\ \citep{datta2013}.

\section{MERW and Hub Identification}
\label{sec:merw}

\subsection{Maximal Entropy Random Walk}

The Maximal Entropy Random Walk \citep{burda2009localization} is the unique stationary Markov chain on $G$ that assigns equal probability to all paths of equal length between any two fixed endpoints. It is the maximally unbiased random walk consistent with the graph structure.

Let $A \in \{0,1\}^{N \times N}$ be the adjacency matrix of $G$. By the Perron--Frobenius theorem, $A$ has a unique positive leading eigenvalue $\lambda_1$ with a corresponding strictly positive eigenvector $\psi \in \mathbb{R}^N_{>0}$, normalized so that $\|\psi\|_2 = 1$. The MERW transition matrix is
\[
S_{ij} = \frac{A_{ij}}{\lambda_1} \cdot \frac{\psi_j}{\psi_i},
\]
and its unique stationary distribution is $\pi_i = \psi_i^2$. Unlike the standard random walk, whose stationary distribution is proportional to degree, MERW concentrates its stationary mass on structurally central nodes. We use $\psi_i$ as the centrality score of node $i$; see Appendix~\ref{app:merw} for a detailed discussion of the localization phenomenon and its information-theoretic justification.

\subsection{Decentralized Hub Identification}

In a multi-agent setting, no single agent has access to the full adjacency matrix $A$. We compute $\psi$ via the distributed power iteration of \citep{jelasity2007asynchronous}, in which each agent normalizes using only local gossip rather than a global reduce. Crucially, the gossip normalization replaces the standard all-reduce with pairwise averaging of log-growth rates, so no agent ever needs to see the global topology. Since we require only the \emph{ratio} $\psi_j / \psi_i$ to determine routing directions, and not the absolute eigenvector values, the unnormalized iterates suffice; see Appendix~\ref{app:poweriter} for the full protocol and convergence analysis.

After power iteration, each agent $i$ holds $\tilde{\psi}_i \approx c \cdot \psi_i$ for some common unknown scale $c > 0$. A short \emph{max-flooding} phase propagates the global maximum across the graph: each agent relays the largest $\tilde{\psi}$ value it has seen, and every local maximum re-routes toward the incoming direction. The result is a single directed spanning tree $\mathcal{T}$ rooted at the unique global maximum $h^\star = \arg\max_i \tilde{\psi}_i$, with all edges oriented from children toward the root. The phase runs for a fixed hyperparameter number of rounds $T_{\mathrm{flood}} \geq \mathrm{diam}(G)+1$, enough for the maximum to reach every node; see Appendix~\ref{app:flooding} for the detailed protocol.

\begin{proposition}[Centralized aggregation via routing]
\label{prop:aggregation}
Suppose each node $i \in V$ holds a scalar $x_i$ and propagates it toward hub $h^\star$ one hop per round, with each intermediate node forwarding the sum of what it receives. After $D$ rounds, the hub holds $\sum_{i \in V} x_i$.
\end{proposition}

\begin{proof}
By induction on tree depth. Leaf nodes (depth $D$) send $x_i$ to their parents after round 1. At each subsequent round, every node at depth $d$ receives the partial sums from all its children and forwards the total to its parent. After $D$ rounds the hub has accumulated all contributions.
\end{proof}

This aggregation property is the bridge between the decentralized routing structure and centralized statistical performance.

\begin{proposition}[Broadcast completeness]
\label{prop:broadcast}
Suppose the hub $h^\star$ holds a value $v$ and sends it one hop per round toward its children, with each intermediate node forwarding the received value to all its children. After $D$ rounds, every node $i \in V$ holds $v$.
\end{proposition}

\begin{proof}
By induction on depth. Nodes at depth 1 (direct children of $h^\star$) receive $v$ after round 1. If every node at depth $d-1$ holds $v$ after $d-1$ rounds, then in round $d$ each of them forwards $v$ to all their children (depth $d$). After $D$ rounds all nodes at depth $D$ have received $v$, completing the induction.
\end{proof}

\begin{theorem}[Coordinator sufficiency]
\label{thm:coordinator}
Let $\mathcal{A}$ be any bandit algorithm whose decision at each step depends only on the global sufficient statistics $\{(\hat{n}^k, \hat{s}^k)\}_{k=1}^K$ and the total pull count $P = \sum_{i \in V}\sum_{k=1}^K n_i^k$, where $\hat{n}^k = \sum_{i \in V} n_i^k$ and $\hat{s}^k = \sum_{i \in V} s_i^k$. Suppose a coordinator $h^\star$ exists that:
\begin{enumerate}
    \item receives exact global sufficient statistics from all agents within $D$ rounds (Proposition~\ref{prop:aggregation}), and
    \item computes $P$ by counting the uplink messages received each cycle plus its own pull, and broadcasts $(\hat{n}, \hat{s}, P)$ back to all agents within $D$ rounds (Proposition~\ref{prop:broadcast}), and
    \item all agents begin the next decision round simultaneously (synchronization via the depth signal $D$).
\end{enumerate}
Then the coordinator's information state at each decision point is identical to that of a centralized instance of $\mathcal{A}$. The coordinator therefore achieves the same statistical performance as $\mathcal{A}$ run on the pooled observations of all $N$ agents, with an additive communication overhead of at most $2D$ time steps per decision round.
\end{theorem}

\begin{proof}
Fix any decision round $r$. By Proposition~\ref{prop:aggregation}, after $D$ uplink steps the hub holds $\hat{n}^k = \sum_i n_i^k$ and $\hat{s}^k = \sum_i s_i^k$ for every arm $k$. The hub also counts the number of uplink messages received and adds one for its own pull to obtain $P = \sum_i \sum_k n_i^k$, the exact total number of arm pulls made by all agents in all rounds up to $r$. Together $(\hat{n}, \hat{s}, P)$ are exactly the sufficient statistics a centralized agent would maintain after observing the same pulls. Because $\mathcal{A}$'s decision depends only on these statistics, the hub computes the same action as the centralized agent. By Proposition~\ref{prop:broadcast}, after $D$ downlink steps every agent holds $(\hat{n}, \hat{s}, P)$ and the hub's decision. By the depth signal, all agents begin the next pull phase at the same time step. Therefore, in terms of arm pulls, the protocol is a faithful simulation of $\mathcal{A}$ on the pooled observations, and any performance guarantee for $\mathcal{A}$ carries over. The only cost is $2D$ idle time steps per round for communication.
\end{proof}

Figures~\ref{fig:merw_viz_ba}--\ref{fig:merw_viz_grid} illustrate the MERW centrality and induced routing tree on three representative graph families. On the Barab\'{a}si--Albert graph, $\psi^2$ concentrates sharply on the highest-degree hub. On the Erd\H{o}s--R\'{e}nyi graph, centrality is more spread but still identifies a clear hub. On the grid, MERW correctly selects the center node. Additional graph families (barbell, lollipop, chain, star, clustered) are shown in Appendix~\ref{app:merw_viz}.

\begin{figure}[H]
    \centering
    \includegraphics[width=\textwidth]{results/MERW_visualization/merw_tree_ba_N20.pdf}
    \caption{MERW centrality and routing tree on a Barab\'{a}si--Albert graph ($N=20$). Three panels: \emph{left} -- undirected graph with nodes colored by $\psi_i/\psi_{\max}$; \emph{center} -- routing subtree overlaid on the graph, hub ringed in gold; \emph{right} -- induced directed tree in hierarchical layout, hub at top, arrows pointing child $\to$ parent.}
    \label{fig:merw_viz_ba}
\end{figure}

\begin{figure}[H]
    \centering
    \includegraphics[width=\textwidth]{results/MERW_visualization/merw_tree_er_N20.pdf}
    \caption{MERW centrality and routing tree on an Erd\H{o}s--R\'{e}nyi graph ($N=20$). Panel layout as in Figure~\ref{fig:merw_viz_ba}.}
    \label{fig:merw_viz_er}
\end{figure}

\begin{figure}[H]
    \centering
    \includegraphics[width=\textwidth]{results/MERW_visualization/merw_tree_grid_N25.pdf}
    \caption{MERW centrality and routing tree on a grid graph ($N=25$). Panel layout as in Figure~\ref{fig:merw_viz_ba}.}
    \label{fig:merw_viz_grid}
\end{figure}

\subsection{The Two Versions of \merw}

The hub and tree are constructed identically for both the regret and BAI objectives. We describe each version in the following sections.

\section{Regret Minimization}
\label{sec:regret}

\subsection{Protocol}

There are $K$ arms with unknown means $\mu_1 \geq \mu_2 \geq \cdots \geq \mu_K$ and sub-Gaussian reward distributions with parameter $\sigma^2$. The objective is to minimize the expected cumulative group regret:
\[
R(T) = \sum_{i \in V} \sum_{t=1}^{T} \mathbb{E}[\mu_1 - \mu_{a_i(t)}].
\]

Each relay cycle lasts exactly $2D + 1$ rounds. At the start of each cycle all $N$ agents pull simultaneously; the hub aggregates observations via the uplink, runs UCB on the global statistics, and broadcasts its estimates back down. Non-hub nodes update their local estimates upon receiving the downlink and schedule their next pull to align with the start of the next cycle. The total-pull count $P$ (indexed by actual arm pulls, not wall-clock rounds, following \citep{martinez2019decentralized}) is maintained exactly by the hub and broadcast in the downlink so all nodes use a consistent UCB bonus. \merw-UCB is one instantiation of the generic protocol given as Algorithm~\ref{alg:merw-generic} (Appendix~\ref{app:pseudocode}); Table~\ref{tab:instantiations} in the same appendix specifies the hooks that turn Algorithm~\ref{alg:merw-generic} into \merw-UCB and, side by side, into \merw-BAI (Section~\ref{sec:bai}).

\subsection{Regret Bound}

\begin{corollary}[Regret matches centralized UCB]
\label{cor:regret}
Under the \merw{} relay protocol, the expected group cumulative regret satisfies \textcolor{red}{[PLACEHOLDER: regret bound to be stated here]}, where the leading term matches a centralized UCB agent on $P$ total pulls and the dependence on wall-clock time $T$ enters through $P = \Theta(NT/(2D+1))$.
\end{corollary}

\begin{proof}[Proof sketch]
By Theorem~\ref{thm:coordinator}, the hub's information state $(\hat{n}^k, \hat{s}^k, P)$ is identical to that of a centralized UCB agent at every decision point, so the standard UCB regret analysis applies directly to the hub's pull sequence indexed by $P$. During the $2D$ communication steps of each cycle, only the hub pulls once; non-hub agents are idle and accumulate no regret. The dominant cost is therefore not an additive penalty per cycle but a reduced pull rate: each cycle of $2D+1$ wall-clock steps produces exactly one pull per agent, so over horizon $T$ the total pulls are $P = \Theta(NT/(2D+1))$. The regret bound in terms of $T$ follows by substituting this expression for $P$ into the centralized UCB guarantee.
\end{proof}

\subsection{Empirical Results}

We evaluate \merw{} on Barab\'{a}si--Albert and Erd\H{o}s--R\'{e}nyi graphs ($N=20$, $K=5$, $T=5000$, 50 runs, Figures~\ref{fig:regret_ba}--\ref{fig:regret_er}). Results on additional graph families are given in Appendix~\ref{app:regret}.

Each figure shows three panels. The first panel shows group cumulative regret $\sum_i R_i(t)$ as a function of wall-clock round $t$; communication-only rounds are included, so the curve is flat during uplink and downlink phases and steps up only when pulls occur. The second panel shows the hub's individual regret $R_{h^\star}(t)$, isolating the performance of the emergent coordinator. The third panel shows group regret as a function of total arm pulls across all agents, removing the communication overhead and measuring pure sample efficiency.

On both graph families the hub, identified from local MERW computations without any central authority, successfully coordinates the network: group regret grows logarithmically and the hub's individual regret tracks the same curve, confirming that the emergent coordinator behaves as a centralized UCB agent. The regret-vs.-pulls curve shows that sample efficiency is not degraded by the relay structure. Results on additional graph families (barbell, lollipop, chain, star, grid) are given in Appendix~\ref{app:regret}.

\begin{figure}[H]
    \centering
    \includegraphics[width=\textwidth]{results/Regret/relay_regret_ba_N20_K5_T5000.pdf}
    \caption{Cumulative regret on a Barab\'{a}si--Albert graph ($N=20$, $K=5$, $T=5000$, 50 runs). \emph{Left}: group regret vs.\ round; \emph{center}: hub regret vs.\ round; \emph{right}: group regret vs.\ total arm pulls. Shaded bands are standard error of the mean.}
    \label{fig:regret_ba}
\end{figure}

\begin{figure}[H]
    \centering
    \includegraphics[width=\textwidth]{results/Regret/relay_regret_er_N20_K5_T5000.pdf}
    \caption{Cumulative regret on an Erd\H{o}s--R\'{e}nyi graph ($N=20$, $K=5$, $T=5000$, 50 runs). Panel layout as in Figure~\ref{fig:regret_ba}.}
    \label{fig:regret_er}
\end{figure}

\section{Best Arm Identification}
\label{sec:bai}

\subsection{Setting}

In the BAI setting the horizon $T$ is not fixed. Instead, given a confidence parameter $\delta \in (0,1)$, the goal is to identify the best arm $a^\star = \arg\max_k \mu_k$ with probability at least $1-\delta$ while minimizing the total number of arm pulls across all agents.

\subsection{\merw-BAI: Hillel Elimination over the MERW Tree}

\merw-BAI adapts Hillel's successive elimination to the tree topology induced by MERW. The pull schedule and elimination rule are identical to Hillel's Algorithm 3; the difference is the communication model. Rather than all-to-all broadcast, agents communicate via the MERW spanning tree: local sums travel hop-by-hop up to the hub (uplink), the hub computes the elimination, and the surviving set $S_r$ together with the timing signal $D$ travels back down (downlink). The timing signal $D$ ensures all nodes restart the next round simultaneously, regardless of their depth in the tree.

\paragraph{Round structure.} Each elimination round $r$ has four phases:

\begin{enumerate}
    \item \textbf{Pull phase} ($L_r$ steps): every node independently computes $L_r = \lceil t_r - t_{r-1} \rceil$ from the pre-agreed schedule (all nodes know $N$, $K$, $\delta$, so no communication is needed). Each node pulls every arm in $S_{r-1}$ exactly $L_r$ times and accumulates local sums.

    \item \textbf{Uplink} ($D$ steps): local sums propagate hop-by-hop toward the hub, one hop per time step. A node at depth $d$ reaches the hub after $d$ steps.

    \item \textbf{Hub aggregation}: the hub sums all received packets to form the global aggregate $\hat{s}_k = \sum_{i=1}^N s_k^{(i)}$ for each surviving arm $k$. It then computes global averages $\tilde{p}_k = \hat{s}_k / (N \cdot L_r)$, identifies $\tilde{p}^\star = \max_k \tilde{p}_k$, and eliminates arms where $\tilde{p}_k < \tilde{p}^\star - \epsilon_r$.

    \item \textbf{Downlink} ($D$ steps): the hub broadcasts the new surviving set $S_r$ and the depth signal $D$ hop-by-hop back down the tree. Each node at depth $d$ receives the packet after $d$ steps and waits $D - d$ additional steps before starting the next pull phase, ensuring all nodes are synchronized.
\end{enumerate}

The total time cost per round is $L_r + 2D$ steps (pull + uplink + downlink), compared to $L_r + 1$ for Hillel (pull phase plus one all-to-all broadcast round). The sample complexity in terms of total arm pulls is identical: $N \cdot L_r \cdot |S_{r-1}|$ per round. \merw-BAI is the second instantiation of the generic protocol, Algorithm~\ref{alg:merw-generic} (Appendix~\ref{app:pseudocode}); Table~\ref{tab:instantiations} specifies its hooks side by side with \merw-UCB's.

\subsection{Comparison with Hillel}

The two algorithms are identical in sample complexity by construction: both use the same $L_r$ schedule and the same elimination rule. The difference lies entirely in the communication model:

\begin{itemize}
    \item \textbf{Hillel}: assumes all-to-all broadcast in one round. The coordinator is not needed because every agent computes the elimination independently. Requires $O(N)$ messages per phase.
    \item \textbf{\merw-BAI}: uses the MERW tree. The hub acts as the aggregation point because it is the only node with the full global picture. Requires $O(N)$ messages per phase (each node sends one packet), but communication takes $O(D)$ time steps rather than $O(1)$.
\end{itemize}

The central claim is that the hub, which plays the role of aggregator in \merw-BAI, is \emph{identified from the graph structure itself} via decentralized MERW computation; no all-to-all communication is assumed. This makes \merw-BAI applicable to any undirected graph, at the cost of $O(D)$ extra time steps per round relative to Hillel. The sample complexity equivalence is formalized in Corollary~\ref{cor:bai}.

\begin{corollary}[BAI sample complexity matches Hillel]
\label{cor:bai}
Under the \merw-BAI protocol, with $\mathcal{A} = $ Hillel's successive elimination (Algorithm 3 of \citep{hillel2013distributed}), the total number of arm pulls to identify the best arm with probability at least $1-\delta$ is identical to that of Hillel's algorithm. The only overhead is $2D$ extra time steps per elimination round.
\end{corollary}

\begin{proof}
By Theorem~\ref{thm:coordinator}, the hub receives exact global sums $\hat{s}_k = \sum_i s_k^{(i)}$ after $D$ uplink steps, computes global averages $\tilde{p}_k = \hat{s}_k/(N \cdot L_r)$, and applies the elimination threshold $\epsilon_r$ identically to Hillel's centralized aggregator. The surviving set $S_r$ is therefore the same as in Hillel's algorithm. By Proposition~\ref{prop:broadcast} and the depth signal, $S_r$ reaches every agent within $D$ downlink steps, and all agents enter round $r+1$ simultaneously. Since $L_r$ and $\epsilon_r$ are pre-agreed and identical across agents, the total pull count $N \cdot \sum_r L_r \cdot |S_{r-1}|$ matches Hillel's exactly. The correctness guarantee $\Pr[\text{output} = a^\star] \geq 1-\delta$ follows from Hillel's analysis applied to the hub's elimination sequence.
\end{proof}

Empirically, we compare total arm pulls as a function of $N$ in Figure~\ref{fig:bai}, where both algorithms use the same graph and the same hub node (identified by MERW) as the aggregation point for a fair comparison.

\begin{figure}[H]
    \centering
    \includegraphics[width=0.85\textwidth]{results/BAI/merw_ucb_bai_ba_K10.pdf}
    \caption{BAI sample complexity on Barab\'{a}si--Albert graphs ($K=10$, $\delta=0.05$, 50 runs per $N$, fixed graph topology per $N$). \emph{Left}: total arm pulls vs.\ number of agents $N$. \emph{Right}: pulls per agent vs.\ $N$. Error bars are standard error of the mean. Both algorithms use the MERW-identified hub as aggregation point; \merw-BAI routes communication through the induced spanning tree while Hillel assumes direct all-to-all broadcast.}
    \label{fig:bai}
\end{figure}

\section{Fault Tolerance}
\label{sec:ft}

Every node in $\mathcal{T}$ is a potential failure point, and \merwft{} must account for the failure of \emph{any} node, not only the hub. The consequences, however, are highly asymmetric. The hub is the coordinator: it alone runs the bandit algorithm on globally aggregated statistics, and if it fails permanently, no downlink is broadcast, no agent learns the outcome of the current round, and the entire coordination structure collapses -- restarting the full \merw{} pipeline mid-episode is infeasible, since it requires a fresh power-iteration and max-flooding phase with no accumulated statistics. A non-hub node, by contrast, is a relay for one subtree: its failure disrupts only the agents below it, while the hub and the rest of the tree continue operating normally. Both cases must be handled, but they call for different mechanisms and different urgency, so we treat them in turn: Section~\ref{sec:ft-backup} through Section~\ref{sec:ft-recursive} address hub failure, the more consequential case, and Section~\ref{sec:nonhub} addresses non-hub node failure, which is more localized and, as we show, cheaper to repair.

\merwft{} tolerates permanent hub failure with zero communication overhead during normal operation. The key idea is to designate the hub's direct child with the highest eigenvector centrality as a passive backup. The backup mirrors the hub's global state at every cycle for free, since it already receives the full downlink broadcast. Upon detecting hub silence, the backup promotes itself as the new root, re-routes its sibling nodes at depth 1, and resumes the bandit protocol from the last known state. Failover completes in $O(\mathrm{diam}(G))$ rounds in general and $O(1)$ rounds when the backup has direct graph edges to all former hub siblings. The backup's state is at most one cycle stale at the moment of promotion. Non-hub failure, addressed in Section~\ref{sec:nonhub}, is repaired locally via precomputed swap edges in $O(1)$ rounds, without ever involving the hub.

\subsection{Background and Setup}
\label{sec:ft-setup}

We assume the same setting as \merw: $N$ agents on an undirected connected graph $G = (V, E)$, a spanning tree $\mathcal{T}$ rooted at hub $h^\star$ constructed via \merw{} power iteration and max-flooding. Immediately after max-flooding, each node $i$ knows only its own parent $\mathrm{par}(i)$ (its own routing decision); it does not yet know its children $\mathrm{ch}(i)$, its depth $d_i$, or the tree height $D$. These are discovered over the course of the first cycle: $i$ learns $\mathrm{ch}(i)$ the first time each child routes a packet through it, learns $d_i$ from the hop count carried on the first downlink it receives, and learns $D$ only once the hub has broadcast it -- the hub itself only knows $D$ after observing that its uplink has gone silent for a round. Each relay cycle consists of an uplink phase ($D$ steps), a hub aggregation step, and a downlink phase ($D$ steps). At the end of each downlink, every node at depth 1 holds the hub's global state $(\hat{n}, \hat{s}, P)$.

We add two assumptions not present in the base \merw{} protocol.

\begin{definition}[Fail-stop failure]
    A node may fail permanently at an unknown round $t_f$. After $t_f$, it sends no
    messages and processes no packets. All other nodes are assumed to be live. Failure
    is detectable via timeout: if a node expects a message within a known window and
    receives none, it concludes the sender has failed.
\end{definition}

The timeout-based failure detector is necessary. Anagnostou and Hadzilacos
\citep{anagnostou1993} prove that no purely self-stabilizing protocol can distinguish a
permanently crashed node from a very slow live node in an asynchronous network; an
external failure oracle (here, the cycle-length timeout) is required. In the fully
asynchronous model, Chandra and Toueg \citep{chandratoueg1996} show that even
timeout-based detectors are inherently unreliable: they can only suspect a node, not
certify its failure. In our setting this ambiguity is resolved by the synchronous round
structure of the bandit protocol: a missing uplink packet at a round boundary is
conclusive evidence of failure, not delay.

\begin{definition}[2-connectivity]
\label{def:2conn}
    The underlying graph $G$ is \emph{2-connected}: it remains connected after the
    removal of any single node. Equivalently, $G$ has no cut vertices.
\end{definition}

This assumption is essential for rerouting after any node failure. The routing tree
$\mathcal{T}$ is a spanning tree of $G$, not $G$ itself. During normal operation only
tree edges are used; non-tree edges $E \setminus E(\mathcal{T})$ are idle. When a node
$v$ fails, $\mathcal{T}$ loses $v$ and becomes a forest, but $G \setminus \{v\}$ remains
connected by 2-connectivity. Rerouting therefore always uses non-tree edges of $G$ to
re-establish paths that were broken in $\mathcal{T}$. Each node knows its neighbors in
$G$ (not just its parent and children in $\mathcal{T}$), since neighbors in $G$ are
discovered during the power iteration neighbor-exchange phase.

\begin{remark}
    If $G$ is a tree itself (i.e., $\mathcal{T} = G$), then $G$ has no non-tree edges
    and every node is a cut vertex. Any single node failure permanently disconnects some
    part of the network and no rerouting is possible. The 2-connectivity assumption is
    therefore tight: it is the minimal condition that guarantees rerouting always succeeds
    after a single node failure.
\end{remark}

\subsection{The Backup Hub}
\label{sec:ft-backup}

\subsubsection{Selection}

During max-flooding, the hub $h^\star$ is the global maximum and never receives a flood
packet from any neighbor: the flood propagates toward $h^\star$, not outward from it.
Nodes route uplink packets toward the hub silently, without announcing their routing
decision. The hub therefore does not know $\mathrm{ch}(h^\star)$ until the first uplink
of the bandit phase, when packets from depth-1 nodes arrive directly.

The hub learns its children from the first uplink: every node $j$ at depth 1 sends its
observations directly to $h^\star$, revealing its identity. Along with each uplink packet,
$j$ includes its local $\tilde{\psi}_j$ value (already computed during power iteration
and stored locally). The hub collects all arriving packets, identifies
$\mathrm{ch}(h^\star)$, and selects the backup:
\[
    b^\star = \arg\max_{j \in \mathrm{ch}(h^\star)} \tilde{\psi}_j.
\]
The hub announces $b^\star$ in the \emph{first downlink}, adding a single identifier
field. All agents store $b^\star$'s identity upon receiving it. The first cycle therefore
runs without a backup; hub failure during this single cycle is not covered by \merwft.

\subsubsection{Analogy with EIGRP DUAL}
\label{sec:eigrp}

The selection of $b^\star$ is directly analogous to the Diffusing Update Algorithm (DUAL)
used in EIGRP \citep{garcia1993eigrp}, the routing protocol deployed in Cisco networks.
We describe DUAL briefly to make the analogy precise.

\paragraph{DUAL setup.} Each router $i$ maintains a routing table entry for every
destination $d$, consisting of:
\begin{itemize}
    \item $\mathrm{FD}_i(d)$: the \emph{feasible distance}, the best known distance from
        $i$ to $d$ at the time the current successor was chosen.
    \item $\mathrm{RD}_j(d)$: the \emph{reported distance} from neighbor $j$ to $d$,
        as advertised by $j$.
    \item $\mathrm{successor}(i, d)$: the next-hop neighbor currently used to reach $d$.
    \item $\mathrm{FS}(i, d)$: the set of \emph{feasible successors}, pre-computed
        backup next-hops.
\end{itemize}

\paragraph{Feasibility condition.} A neighbor $j$ is a feasible successor for destination
$d$ at router $i$ if:
\[
    \mathrm{RD}_j(d) < \mathrm{FD}_i(d).
\]
This condition guarantees that $j$ is strictly closer to $d$ than $i$ was when $i$ last
updated its route, which is sufficient to ensure loop-freedom: $j$ cannot be routing
through $i$ to reach $d$.

\paragraph{Failover.} When the current successor to $d$ fails, DUAL checks whether
$\mathrm{FS}(i, d)$ is non-empty. If so, the feasible successor with the lowest
$\mathrm{RD}_j(d)$ is promoted to successor immediately, with no query to the rest of
the network. Convergence is sub-second. If no feasible successor exists, DUAL initiates
a diffusing computation (a network-wide query) to find a new route, which takes $O(D)$
rounds in the worst case.

\paragraph{Analogy with \merwft.} In our setting, the ``destination'' is the hub
$h^\star$ and the ``distance'' metric is replaced by eigenvector centrality $\tilde{\psi}$.
The feasibility condition becomes: $b^\star$ is a valid backup for $h^\star$ if
$\tilde{\psi}_{b^\star}$ is the highest among $h^\star$'s direct children, i.e.,
$b^\star$ is the most structurally central node adjacent to $h^\star$ in the tree. This
is loop-free by construction: $b^\star$ is a child of $h^\star$, so it routes toward
$h^\star$ and cannot already be routing through $h^\star$ to reach itself. The promotion
in Section~\ref{sec:promotion} is the \merw{} analog of DUAL's feasible-successor activation:
immediate, local, and requiring no network-wide query. When no backup exists (hub has
only one child, or all children have equal $\tilde{\psi}$), the protocol falls back to a
full \merw{} re-run, analogous to DUAL's diffusing computation.

\subsubsection{Passive State Mirroring}

During normal operation, $b^\star$ receives every downlink broadcast from $h^\star$
exactly as any other depth-1 node does. It stores the received state:
\[
    (\hat{n}^{b^\star}, \hat{s}^{b^\star}, P^{b^\star}) \leftarrow (\hat{n}, \hat{s}, P)
\]
at the end of each cycle. No extra messages are sent; no extra computation is performed.
This is passive replication in the sense of Lamport, Malkhi and Zhou
\citep{lamport2009vertical}: the backup mirrors primary state synchronously and for free,
since it already lies on the downlink path.

\begin{proposition}[Backup state staleness]
\label{prop:staleness}
    At the end of cycle $r$, the backup $b^\star$ holds $(\hat{n}, \hat{s}, P)$ equal to
    the hub's state after the hub's own pull in cycle $r$. If the hub fails during cycle
    $r+1$ at any point before completing its downlink, the backup's state reflects all
    observations up to and including the hub's pull in cycle $r$.
\end{proposition}

\begin{proof}
    In cycle $r$, the hub receives the full uplink, updates $(\hat{n}, \hat{s})$ with all
    agents' observations, increments $P$, performs its own pull, and broadcasts the
    updated state in the downlink. $b^\star$ receives this downlink at step $d_{b^\star} =
    1$ (it is at depth 1), so it holds the hub's state after the hub's pull. If the hub
    fails at any point in cycle $r+1$ before completing the downlink of cycle $r+1$,
    the last downlink $b^\star$ received was from cycle $r$. The backup's state is
    therefore at most one full cycle stale.
\end{proof}

\subsection{Failover Protocol}

\subsubsection{Failure Detection}
\label{sec:failure-detection}

Each depth-1 node, including $b^\star$, expects a downlink packet from $h^\star$ within
$D$ steps of the end of the uplink phase. If no packet arrives, it declares the hub
failed. Since all depth-1 nodes are synchronized (they sent their uplink packets at the
same wall-clock step), they all detect silence simultaneously after the same timeout.

\subsubsection{Promotion}
\label{sec:promotion}

Upon detecting hub failure, $b^\star$ executes the following promotion procedure:

\begin{enumerate}
    \item \textbf{Self-promotion}: $b^\star$ sets $d_{b^\star} \leftarrow 0$,
        $\mathrm{par}(b^\star) \leftarrow \text{undefined}$, and begins acting as the
        hub using its stored state $(\hat{n}^{b^\star}, \hat{s}^{b^\star}, P^{b^\star})$.

    \item \textbf{Sibling re-routing}: $b^\star$ must notify all former siblings
        $\mathrm{ch}(h^\star) \setminus \{b^\star\}$ that it is the new hub. In
        $\mathcal{T}$, $b^\star$ is not directly connected to its siblings: their only
        shared parent was $h^\star$, now failed. However, by 2-connectivity of $G$
        (Definition~\ref{def:2conn}), each sibling $s$ has at least one neighbor in
        $G \setminus \{h^\star\}$ other than its own subtree. $b^\star$ floods the
        promotion message $\langle b^\star, \text{NEW-HUB}, d=1 \rangle$ over $G$-edges:
        it sends to all its $G$-neighbors, which forward to their $G$-neighbors, until
        every former sibling is reached. Since $G \setminus \{h^\star\}$ is connected
        (by 2-connectivity), the flood reaches all siblings in at most $\mathrm{diam}(G
        \setminus \{h^\star\})$ steps. Each sibling $s$ upon receiving the message sets
        $\mathrm{par}(s) \leftarrow b^\star$ and $d_s \leftarrow 1$ if it has a direct
        $G$-edge to $b^\star$, or re-routes via an intermediate node otherwise. Subtrees
        below depth 1 are unaffected, subject to the synchronization mechanism of
        Section~\ref{sec:sync-during-reroute}.

    \item \textbf{Depth update}: $b^\star$ computes the new tree height. Since $b^\star$
        moves from depth 1 to depth 0, its own subtree's depths decrease by 1. However,
        the former siblings $s \in \mathrm{ch}(h^\star) \setminus \{b^\star\}$ now route
        through $b^\star$, so nodes in their subtrees gain one extra hop. The new height
        is:
        \[
            D' = \max\!\bigl(D - 1,\; \max_{s \in \mathrm{ch}(h^\star) \setminus \{b^\star\}} (d_s^{\max} + 1)\bigr),
        \]
        where $d_s^{\max}$ is the depth of the deepest leaf in $s$'s subtree. In the
        common case where all depth-1 subtrees have the same height, $D' = D$. In the
        worst case (e.g., a star graph where the hub has only one child $b^\star$ with no
        subtree and all other agents are direct children of $h^\star$), the siblings
        become children of $b^\star$ at depth 1, giving $D' = 1 = D$, so no increase
        occurs. $b^\star$ broadcasts $D'$ in the first downlink after promotion so all
        agents synchronize on the new cycle length.

    \item \textbf{Resume}: $b^\star$ immediately begins the next cycle as the new hub,
        accepting uplink packets from its children (now including the re-routed siblings).
\end{enumerate}

\begin{proposition}[Failover correctness]
\label{prop:failover}
    After promotion, the network forms a valid spanning tree $\mathcal{T}'$ rooted at
    $b^\star$ with height $D' \leq D + 1$. Every node has a unique path to $b^\star$,
    and no cycles are introduced.
\end{proposition}

\begin{proof}
    Before promotion, $\mathcal{T}$ is a directed tree rooted at $h^\star$. Remove
    $h^\star$ from $\mathcal{T}$. The result is a forest: one subtree rooted at each
    child of $h^\star$, including $b^\star$. Step 2 adds a directed edge from each
    sibling $s \in \mathrm{ch}(h^\star) \setminus \{b^\star\}$ to $b^\star$, merging
    all subtrees into a single tree rooted at $b^\star$. Since the original edges within
    each subtree are unchanged and each subtree was cycle-free, and the new edges are
    directed from former siblings to $b^\star$ (which has no parent), no cycle can form.
    The height increases by at most 1 (the path from the deepest node in a former sibling
    subtree now passes through $b^\star$ rather than $h^\star$, which is the same number
    of hops). Every node has a unique directed path to $b^\star$ by construction.
\end{proof}

\begin{proposition}[Failover completeness]
\label{prop:completeness}
    The promotion procedure completes in $O(\mathrm{diam}(G))$ rounds and requires
    $O(|E|)$ messages in the worst case. If $b^\star$ is directly connected in $G$ to
    all its former siblings, it completes in $O(1)$ rounds and $O(|\mathrm{ch}(h^\star)|)$
    messages.
\end{proposition}

\begin{proof}
    Step 1 is local. Step 2 floods the promotion message over $G \setminus \{h^\star\}$,
    which is connected by 2-connectivity; the flood reaches all nodes in at most
    $\mathrm{diam}(G \setminus \{h^\star\}) \leq \mathrm{diam}(G)$ rounds, touching at
    most $|E|$ edges. If every former sibling has a direct $G$-edge to $b^\star$, step 2
    completes in 1 round with $|\mathrm{ch}(h^\star)| - 1$ messages. Steps 3 and 4 are
    piggy-backed on step 2 at no extra cost.
\end{proof}

\subsection{Non-Hub Node Failure}
\label{sec:nonhub}

When a non-hub node $v$ at depth $d_v$ fails, the impact is local: only the subtree
rooted at $v$ is disrupted. The hub and all nodes outside this subtree continue operating
normally.

\paragraph{Failure detection.} The correct locus of detection is $v$'s parent
$\mathrm{par}(v)$, not $v$'s children and not the hub. In each round, $\mathrm{par}(v)$
expects a forwarded uplink packet from $v$ within a known window (one hop delay). If no
packet arrives, $\mathrm{par}(v)$ declares $v$ failed. This is strictly faster than
child-driven detection (children would only notice failure when their own uplink finds no
relay, which takes $D - d_v$ hops to manifest) and faster than hub-driven detection
(which requires a full cycle to notice a missing contribution). The parent-driven approach
is supported by Chitnis et al.\ \citep{chitnis2009} and the proactive multicast recovery
of Liao et al.\ \citep{liao2007proactive}, who show that parent monitoring minimizes both
detection latency and false-positive rate compared to child or hub detection.

Once $\mathrm{par}(v)$ detects failure, it notifies $v$'s children
$\mathrm{ch}(v)$ directly via $G$-edges (it knows $\mathrm{ch}(v)$ because $v$'s
children sent uplink packets through $v$ to $\mathrm{par}(v)$ in prior rounds, revealing
their identities). It pushes the re-attachment assignment to each child in one message,
eliminating the need for children to independently scan their neighborhood.

\paragraph{Precomputed swap edges.} Rather than scanning $G$-neighbors at failure time,
each node $c$ precomputes a \emph{swap parent} $u_c$ once depths are established,
following the best-swap-edge framework of Nardelli et al.\ \citep{nardelli2003} and its
distributed variant by Datta et al.\ \citep{datta2013}. Depths are not known immediately
after max-flooding (Section~\ref{sec:ft-setup}): each node learns its own $d_i$ only from
the hop count on the first downlink it receives, during the first bandit cycle. Once
every node knows its depth, the precomputation runs as a one-shot side computation
before the next cycle begins: each node broadcasts its depth to its $G$-neighbors, and
the swap parent minimizes the maximum depth increase across $c$'s entire subtree, not
just $c$'s own depth:
\[
    u_c = \arg\min_{u \in \mathcal{N}_G(c) \setminus \{v\}} \max_{\ell \in \mathrm{subtree}(c)} (d_u + 1 + d_\ell - d_c),
\]
where $d_\ell - d_c$ is the depth of $\ell$ within $c$'s subtree. This precomputation
requires one extra message per node (each node broadcasts its depth to its
$G$-neighbors once), costs $O(|E|)$ total messages, and
eliminates any on-the-fly neighbor scanning at failure time. Re-attachment then takes
exactly 1 round: $\mathrm{par}(v)$ pushes the swap assignments, each child $c$ sets
$\mathrm{par}(c) \leftarrow u_c$ and updates its depth as
\[
    d_c' = d_{u_c} + 1,
\]
and propagates the shift to every descendant $\ell$ in its subtree:
\[
    d_\ell' = d_c' + (d_\ell - d_c).
\]
The whole subtree of $c$ shifts uniformly by the offset $d_{u_c} + 1 - d_c$, so the
relative structure within the subtree is preserved.

\paragraph{Depth increase.} If $v$ was the only node at depth $d_v$ on the path between
some leaf $\ell$ and the hub, re-attachment may increase the depth of $\ell$'s subtree.
In the worst case, if $c$ has no neighbor at depth $d_v - 1$ other than $v$, it must
attach to a node at depth $d_v$, making its own depth $d_v + 1$ and pushing all
descendants down by one level. The new tree depth is
\[
    D' = \max\bigl(D_{\mathrm{rest}},\ \max_{c \in \mathrm{ch}(v)} (d_{u_c} + 1 + d_c^{\max} - d_c)\bigr),
\]
where $D_{\mathrm{rest}}$ is the depth of the unaffected part of the tree and
$d_c^{\max}$ is the depth of the deepest leaf in $c$'s subtree before re-attachment.
The hub learns $D'$ via the next uplink: every node piggybacks its current depth, so
the hub takes the max over all received values. In the extreme case of a chain graph
where $v$ is the unique intermediate node at depth $d_v$, this gives $D' = D + 1$.

\paragraph{Consequence for the protocol.} A depth increase changes the cycle length from
$2D + 1$ to $2D' + 1$. The hub must broadcast the updated $D'$ in the next downlink so
all agents synchronize on the new cycle length. This requires one extra downlink field,
identical to the depth signal already present in the protocol. No other structural change
is needed: the bandit algorithm running at the hub is unaffected, and all live agents
continue routing toward the hub via their updated parent pointers.

\begin{proposition}[Non-hub failover correctness]
\label{prop:nonhub}
    After re-attachment of all orphaned children of a failed node $v$, the remaining live
    nodes form a valid spanning tree rooted at $h^\star$ (or $b^\star$ if hub failover
    has already occurred) with height $D' \leq D + |\mathrm{ch}(v)|$. No cycles are
    introduced.
\end{proposition}

\begin{proof}
    Before failure, $\mathcal{T}$ is a directed tree. Remove $v$: the result is a forest
    with one subtree rooted at each $c \in \mathrm{ch}(v)$ and one subtree containing the
    hub. By 2-connectivity of $G$ (Definition~\ref{def:2conn}), $G \setminus \{v\}$ is
    connected, so each orphaned child $c$ has at least one live $G$-neighbor $u_c \in
    \mathcal{N}_G(c) \setminus \{v\}$ (which may be a tree edge to a sibling subtree or
    a non-tree edge of $G$) with finite depth. Since $G \setminus \{v\}$ is connected and
    the hub is in $G \setminus \{v\}$, $u_c$ lies in the hub's component. We select $u_c$
    with $d_{u_c} < d_c$, which is possible because $u_c$ is not itself an orphan (only
    children of $v$ are orphaned under single node failure) and thus has a well-defined
    depth. Adding the directed edge $c \to u_c$ connects $c$'s subtree to the hub's
    component. No cycle forms: $u_c$ routes toward the hub and not through $c$. Repeating
    for each $c \in \mathrm{ch}(v)$ merges all orphaned subtrees into the hub's component.
    The depth of any node in $c$'s
    subtree increases by at most 1 per re-attachment; since there are $|\mathrm{ch}(v)|$
    orphaned children, $D' \leq D + |\mathrm{ch}(v)|$.
\end{proof}

\begin{remark}
    In practice $|\mathrm{ch}(v)|$ is small (the tree degree of an intermediate node),
    so the depth increase is bounded and modest. On graphs with high connectivity,
    orphaned children typically find a neighbor at the same or lower depth, giving
    $D' = D$ with no cycle length change at all.
\end{remark}

\subsection{Discussion}
\label{sec:ft-discussion}

\subsubsection{Why Depth-1 Nodes and Not the Whole Tree}

A natural alternative is to broadcast $b^\star$'s identity and the new depth signal to
the entire tree after promotion, so all nodes update their routing tables. This is
unnecessary. Nodes at depth $\geq 2$ already route toward their parent, which routes
toward $b^\star$ via the unchanged sub-tree structure. The only nodes whose parent
pointer changes are the hub's direct siblings, which are exactly the nodes contacted in
step 2. The rest of the tree is oblivious to the hub failure -- provided, as we show
next, that the re-routing siblings suppress false timeouts in their own subtrees while
step 2 is in flight.

\subsubsection{Synchronizing the Rest of the Tree During Re-Routing}
\label{sec:sync-during-reroute}

Step 2 of promotion can take up to $\mathrm{diam}(G \setminus \{h^\star\})$ rounds
(Proposition~\ref{prop:completeness}). During this window, every former sibling $s$
whose subtree has depth $\geq 1$ is itself waiting to be reached by $b^\star$'s flood,
and so has not yet forwarded a downlink to its own children. Left unaddressed, this
creates a false-timeout cascade: $s$'s children expect a downlink within $D$ steps of
their own uplink (Section~\ref{sec:failure-detection}); if $s$ is still unreached when
that timeout elapses, the children conclude that $s$ itself has failed, triggering an
unrelated non-hub failover (Section~\ref{sec:nonhub}) for a node that is, in fact,
alive and merely waiting on its own re-routing. Without correction, a single hub
failure could therefore spuriously trigger failovers throughout every affected
sibling's subtree.

\paragraph{Fix: propagate a hold signal ahead of re-routing.}
We extend the flood of step 2 with a second, tree-scoped broadcast that each reached
sibling relays downward immediately upon receipt, before it has itself finished
re-routing:

\begin{enumerate}
    \item When $b^\star$'s flood message $\langle b^\star, \text{NEW-HUB}, d{=}1\rangle$
        reaches a former sibling $s$ (possibly via intermediate relay nodes, by
        2-connectivity), $s$ immediately forwards a \textsc{Hold} message down its own
        tree edges to $\mathrm{ch}(s)$, \emph{before} completing its own re-attachment
        to $b^\star$. The \textsc{Hold} message carries no state, only a signal that a
        re-routing is underway one level up.
    \item Every node $v$ that receives \textsc{Hold} resets its own downlink timeout and
        immediately re-forwards \textsc{Hold} to $\mathrm{ch}(v)$, propagating the signal
        to the bottom of $s$'s subtree in $d_s^{\max} - d_s$ additional steps (the depth
        of $s$'s subtree below $s$).
    \item A node that has received \textsc{Hold} does not declare its parent failed
        while waiting; instead, it waits for either (a) a fresh downlink carrying the
        updated $(\hat n, \hat s, P, D')$, which arrives once $s$ completes its own
        re-routing and the next cycle's downlink reaches it, or (b) a second, unrelated
        timeout on the \textsc{Hold} signal itself (see below), whichever comes first.
    \item Since $\mathcal{T}$ is a tree, \textsc{Hold} reaches every node in $s$'s
        subtree along tree edges only, requiring no extra graph exploration beyond what
        the subtree already uses for normal downlink forwarding.
\end{enumerate}

\textsc{Hold} is a strict extension of the existing downlink mechanism, not a parallel
protocol: nodes already forward whatever they receive from their parent to their
children every cycle, so treating \textsc{Hold} as a one-bit downlink payload adds no
new communication primitive. The only new state is a per-node flag, ``re-routing in
progress upstream,'' cleared upon the next ordinary downlink.

\paragraph{Bounding the hold window.} \textsc{Hold} itself must expire, or a
permanently failed $s$ (as opposed to one merely mid-reroute) would leave its subtree
waiting forever. We bound the hold window by the same quantity that bounds step 2: a
node holding \textsc{Hold} for more than $\mathrm{diam}(G)$ additional rounds without
receiving a fresh downlink concludes that $s$ itself has failed, and initiates the
non-hub failover of Section~\ref{sec:nonhub} rooted at $s$. This is strictly later
than the un-held timeout would have fired, so it never masks a genuine failure of $s$;
it only delays the (correct) detection of such a failure until after $b^\star$'s own
re-routing has had a chance to complete first.

\begin{proposition}[No false timeouts during re-routing]
\label{prop:noholdfalsepositive}
    Under the \textsc{Hold} propagation rule above, no node in a re-routing sibling's
    subtree declares its parent failed while step 2 of promotion is still in progress,
    provided step 2 completes within $\mathrm{diam}(G)$ rounds of the sibling receiving
    the flood.
\end{proposition}

\begin{proof}
    Let $s$ be a former sibling and $v$ any node in $s$'s subtree at tree-distance
    $k \geq 1$ below $s$. By construction, $s$ forwards \textsc{Hold} to $\mathrm{ch}(s)$
    in the same round it is reached by $b^\star$'s flood, and every subsequent node
    forwards \textsc{Hold} to its children in the next round upon receipt; hence $v$
    receives \textsc{Hold} exactly $k$ rounds after $s$ is reached, which is before
    $v$'s own downlink timeout would otherwise have elapsed (the timeout is calibrated
    to $D$ rounds from $v$'s uplink, and $s$ is reached within
    $\mathrm{diam}(G\setminus\{h^\star\})$ rounds of hub failure, a quantity independent
    of $v$'s position). Once $v$ holds \textsc{Hold}, by the waiting rule it does not
    declare $s$ failed until an additional $\mathrm{diam}(G)$ rounds have elapsed with no
    fresh downlink. Since step 2 completes within $\mathrm{diam}(G)$ rounds by
    assumption, $s$ finishes re-routing and forwards a fresh downlink to $v$ before $v$'s
    held timeout expires, so $v$ never declares $s$ failed.
\end{proof}

This closes the gap left implicit by the claim that ``the rest of the tree is oblivious
to the hub failure'': obliviousness holds for nodes outside every re-routing sibling's
subtree, but nodes inside such a subtree are not oblivious -- they receive exactly one
extra signal (\textsc{Hold}) that tells them to wait rather than to fail over, which is
enough to prevent a single hub failure from cascading into spurious non-hub failovers
throughout the affected subtrees.

\subsubsection{Recursive Backup: Tolerating Sequential Hub Failures}
\label{sec:ft-recursive}

The single-backup limitation dissolves if the promotion procedure is applied
\emph{recursively}. After $b^\star$ promotes itself to hub, it is now a hub with no
pre-designated successor. Rather than pre-computing a chain of backups, the promoted hub
simply re-runs the same selection procedure it observed from $h^\star$: it learns its new
children from the first uplink, selects the highest-$\tilde\psi$ child as the new backup
$b^{\star\star}$, and announces it in the first downlink. The protocol is therefore
self-renewing: every hub, whether original or promoted, establishes its own successor
after one cycle.

\paragraph{Recursive protocol.} Formally, define the backup selection and announcement
procedure as a subroutine $\textsc{SelectBackup}(h)$:
\begin{enumerate}
    \item Hub $h$ collects the first uplink, learning $\mathrm{ch}(h)$ and their
        $\tilde\psi$ values.
    \item $h$ sets $b_h \leftarrow \arg\max_{j \in \mathrm{ch}(h)} \tilde\psi_j$ and
        announces $b_h$ in the first downlink. All agents store $b_h$.
    \item Upon hub failure, $b_h$ runs the failover protocol of
        Section~\ref{sec:promotion}, promotes itself, and immediately calls
        $\textsc{SelectBackup}(b_h)$.
\end{enumerate}
This recursion terminates naturally when the current hub has no children, meaning it is
the last surviving node and there is nothing left to coordinate.

\begin{proposition}[Sequential hub failure tolerance]
\label{prop:recursive}
    Under the recursive backup protocol, \merwft{} tolerates any finite sequence of
    hub failures $h^\star = h_0, h_1, h_2, \ldots$ where each $h_{k+1}$ is the backup
    selected by $h_k$. After each failure, exactly one cycle runs without a backup, and
    the protocol resumes with a newly selected successor.
\end{proposition}

\begin{proof}
    By induction on the number of failures. At step 0, $h_0 = h^\star$ selects $b_{h_0}$
    after the first uplink (Proposition~\ref{prop:staleness} applies). If $h_0$ fails,
    $h_1 = b_{h_0}$ promotes via Proposition~\ref{prop:failover} and calls
    $\textsc{SelectBackup}(h_1)$. The first cycle of $h_1$'s tenure runs without a
    backup; thereafter $h_1$ has announced $b_{h_1}$. If $h_1$ fails, $h_2 = b_{h_1}$
    promotes and the argument repeats. At each step the number of live nodes decreases by
    at least one (the failed hub), so the sequence terminates in at most $N-1$ failures.
\end{proof}

\paragraph{Depth under sequential failures.} Each promotion may change the tree height
as analyzed in Section~\ref{sec:promotion}. After $k$ sequential hub failures the height
is $D^{(k)} \leq D + k$ in the worst case, since each promotion adds at most one level.
In practice the promoted nodes follow the highest-$\tilde\psi$ path downward from
$h^\star$, which on heterogeneous graphs (Barabasi-Albert, Erdos-Renyi) tends to stay
near the dense core of the graph, keeping $D^{(k)}$ close to $D$.

\subsubsection{Remaining Limitations}

\paragraph{Simultaneous failures.}
If the current hub and its designated backup fail simultaneously, or if the backup fails
before the hub, the protocol has no failover mechanism. Under the recursive protocol,
the unprotected window is exactly one cycle (between hub promotion and backup
announcement); a failure of the new hub during this window leaves the network without a
coordinator. This is a fundamental constraint of the fail-stop model with a single
failure detector: distinguishing simultaneous failure from sequential failure requires
either redundant failure detectors or a Byzantine fault model, both of which are out of
scope.

\paragraph{State staleness and regret.}
Proposition~\ref{prop:staleness} shows the backup's state is at most one cycle stale.
The consequence for bandit performance -- how much regret is incurred by resuming from a
one-cycle-old state -- is left for future work. Informally, UCB's confidence intervals
are monotonically widening in $P$, so using a slightly smaller $P$ makes the algorithm
more exploratory rather than less, which is conservative rather than catastrophic.

\subsubsection{Relation to the Literature}

The passive backup pattern of \merwft{} is an instance of primary-backup replication,
formalized by Lamport, Malkhi and Zhou \citep{lamport2009vertical} as a special case of
Paxos. The one-cycle staleness bound (Proposition~\ref{prop:staleness}) corresponds
exactly to the one-synchronization-round staleness of passive replication.

The impossibility result of Anagnostou and Hadzilacos \citep{anagnostou1993} establishes
why the timeout-based failure detector in Section~\ref{sec:ft-setup} is necessary and not
an engineering convenience: distinguishing a crashed hub from a slow hub is impossible
without it in an asynchronous model.

The sibling re-rooting procedure in Proposition~\ref{prop:failover} is a special case
of the tree reconfiguration algorithm of Arora \citep{arora1994}, which handles arbitrary
node failure including root failure in $O(n)$ time. \merwft{} restricts the structural
change to depth-1 nodes, but the sibling notification must flood over $G$-edges since
$b^\star$ and its former siblings share no direct tree edges after $h^\star$ is removed.
The flood completes in $O(\mathrm{diam}(G))$ in general and $O(1)$ when $b^\star$ has
direct $G$-edges to all former siblings -- a condition that holds with high probability on
dense graphs such as Barabasi-Albert networks, where the hub and its neighbors form a
dense clique-like core. The precomputed swap-edge framework of Nardelli et al.\
\citep{nardelli2003} and its distributed linear-time variant \citep{datta2013} provide the
right foundation for the non-hub re-attachment rule, replacing on-the-fly neighbor
scanning with a one-round precomputed assignment.

\section{Discussion}
\label{sec:discussion}

\textcolor{red}{[TODO: limitations and comments. Key points to address:
\begin{itemize}
    \item The protocol assumes a connected, undirected, static graph. Dynamic or directed graphs are not handled (except for fail-stop node failures, addressed by \merwft; Section~\ref{sec:ft}).
    \item $D$ can be large on sparse or path-like graphs, degrading the effective information rate $N/(2D+1)$.
    \item The BAI version requires agents to know $N$ (for the $L_r$ schedule); the regret version does not.
    \item Power iteration convergence depends on the spectral gap $\lambda_2/\lambda_1$; on nearly-regular graphs MERW centrality may not concentrate sharply.
    \item \merwft{} assumes 2-connectivity and fail-stop failures; Byzantine or simultaneous hub+backup failures are out of scope (Section~\ref{sec:ft-discussion}).
\end{itemize}]}

\section{Conclusion}
\label{sec:conclusion}

\textcolor{red}{[TODO: conclusion. Key points: MERW centrality as a tool for emergent coordination, the hub-and-tree backbone, matching centralized performance for both BAI and regret, the $O(D)$ communication overhead as the price of decentralization, and the passive-backup extension that removes the hub as a single point of failure at zero steady-state cost.]}

\bibliographystyle{plain}
\bibliography{references}

\appendix

\section{Additional MERW Visualizations}
\label{app:merw_viz}

\begin{figure}[H]
    \centering
    \includegraphics[width=\textwidth]{results/MERW_visualization/merw_tree_barbell_N20.pdf}
    \caption{Barbell graph ($N=20$). Panel layout as in Figure~\ref{fig:merw_viz_ba}.}
    \label{fig:merw_viz_barbell}
\end{figure}

\begin{figure}[H]
    \centering
    \includegraphics[width=\textwidth]{results/MERW_visualization/merw_tree_lollipop_N20.pdf}
    \caption{Lollipop graph ($N=20$). Panel layout as in Figure~\ref{fig:merw_viz_ba}.}
    \label{fig:merw_viz_lollipop}
\end{figure}

\begin{figure}[H]
    \centering
    \includegraphics[width=\textwidth]{results/MERW_visualization/merw_tree_chain_N10.pdf}
    \caption{Chain graph ($N=10$). Panel layout as in Figure~\ref{fig:merw_viz_ba}.}
    \label{fig:merw_viz_chain}
\end{figure}

\begin{figure}[H]
    \centering
    \includegraphics[width=\textwidth]{results/MERW_visualization/merw_tree_star_N20.pdf}
    \caption{Star graph ($N=20$). Panel layout as in Figure~\ref{fig:merw_viz_ba}.}
    \label{fig:merw_viz_star}
\end{figure}

\begin{figure}[H]
    \centering
    \includegraphics[width=\textwidth]{results/MERW_visualization/merw_tree_clustered_N24.pdf}
    \caption{Clustered graph ($N=24$). Panel layout as in Figure~\ref{fig:merw_viz_ba}.}
    \label{fig:merw_viz_clustered}
\end{figure}

\section{MERW: Construction and Localization}
\label{app:merw}

\subsection{Construction}

The Maximal Entropy Random Walk \citep{burda2009localization} is the unique stationary Markov chain on $G$ that assigns equal probability to all paths of equal length between any two fixed endpoints. It is the maximally unbiased random walk consistent with the graph structure.

Let $A \in \{0,1\}^{N \times N}$ be the adjacency matrix of $G$. By the Perron--Frobenius theorem, $A$ has a unique positive leading eigenvalue $\lambda_1$ with a corresponding strictly positive eigenvector $\psi \in \mathbb{R}^N_{>0}$, normalized so that $\|\psi\|_2 = 1$. The MERW transition matrix is
\[
S_{ij} = \frac{A_{ij}}{\lambda_1} \cdot \frac{\psi_j}{\psi_i},
\]
and its unique stationary distribution is $\pi_i = \psi_i^2$.

Unlike the standard random walk, whose stationary distribution is proportional to degree ($\pi_i^{\mathrm{SRW}} \propto d_i$), MERW concentrates its stationary mass on a small set of structurally central nodes. This localization phenomenon \citep{burda2009localization} makes $\psi_i$ a natural measure of centrality.

\subsection{Localization and Centrality}

The most striking property of MERW is the sharp concentration of $\pi_i = \psi_i^2$ on a small subset of nodes. On graphs with heterogeneous structure (e.g., Barab\'{a}si--Albert networks), the eigenvector $\psi$ is far from uniform: a handful of nodes carry nearly all the stationary mass, while the rest have negligible values. This mirrors the ground state of a quantum particle on the graph, where probability concentrates in the lowest-energy (most path-rich) region.

We use $\psi_i$ as the centrality score of node $i$. This is justified by two properties:
\begin{enumerate}
    \item $\psi_i$ encodes node $i$'s contribution to paths of \emph{all} lengths in the graph, not just its immediate neighborhood as degree does.
    \item $\psi$ emerges from a maximum-entropy derivation, giving it a clean information-theoretic foundation: it is the least biased measure of structural importance consistent with the graph.
\end{enumerate}

\section{Decentralized Power Iteration}
\label{app:poweriter}

In a multi-agent setting, no single agent has access to the full adjacency matrix $A$. We compute $\psi$ via the distributed power iteration of \citep{jelasity2007asynchronous}, in which each agent normalizes using only local gossip rather than a global reduce.

\paragraph{Protocol.}
Each agent $i$ maintains a local scalar $w_i^{(\tau)}$ initialized to $w_i^{(0)} = 1$. At each iteration $\tau = 1, 2, \ldots, \tau_{\mathrm{init}}$:
\begin{enumerate}
    \item Agent $i$ receives $w_j^{(\tau-1)}$ from all neighbors $j \in \mathcal{N}(i)$ and computes the unnormalized update
    \[
        \tilde{w}_i^{(\tau)} = \sum_{j \in \mathcal{N}(i)} w_j^{(\tau-1)} = (Aw^{(\tau-1)})_i.
    \]
    \item Agent $i$ computes its local log-growth rate $r_i^{(\tau)} = \log(\tilde{w}_i^{(\tau)} / w_i^{(\tau-1)})$, which approximates $\log \lambda_1$ once the iteration has converged.
    \item Agents gossip $r_i^{(\tau)}$ by pairwise averaging with random neighbors for $g$ rounds. After gossip, each $r_i$ approximates the geometric mean growth rate $\frac{1}{N}\sum_j \log(\tilde{w}_j^{(\tau)} / w_j^{(\tau-1)}) \approx \log \lambda_1$.
    \item Each agent normalizes locally: $w_i^{(\tau)} = \tilde{w}_i^{(\tau)} / \exp(r_i^{(\tau)})$.
\end{enumerate}
The iteration runs for a fixed number of steps $\tau_{\mathrm{init}}$, a hyperparameter, so all agents finish on the same synchronous round and enter max-flooding together; no agent evaluates a global convergence test. After $\tau_{\mathrm{init}}$ iterations, each agent holds $\tilde{\psi}_i \approx c \cdot \psi_i$ for some common unknown scale $c > 0$. Convergence error decays geometrically as $O\!\left((\lambda_2/\lambda_1)^{\tau_{\mathrm{init}}}\right)$, so $\tau_{\mathrm{init}}$ trades accuracy for cost, and any value large enough for this error to fall below the routing tolerance suffices.

\paragraph{Why gossip normalization is decentralized.}
The standard power iteration requires a global $\ell_2$ norm at each step, which needs an all-reduce over the entire network. The gossip approach replaces this with $g$ rounds of random pairwise averaging of log-growth rates. Each agent only communicates with one random neighbor per gossip round; no agent needs to know $N$ or the global topology. With $g = O(\log N)$ gossip rounds the approximation error is negligible \citep{jelasity2007asynchronous}.

\paragraph{Why we do not need $N$.}
The key observation is that we use $\psi$ only to determine routing directions: agent $i$ routes toward neighbor $j$ if $\tilde{\psi}_j > \tilde{\psi}_i$. This requires only the \emph{ratio} $\psi_j / \psi_i$, not the absolute value of either component. The gossip normalization produces iterates proportional to $\psi$ up to a common scale, which preserves all relative orderings and is sufficient for hub identification without ever computing $N$.

\paragraph{Cost.} Each iteration requires one neighbor-exchange round plus $g$ gossip rounds, all exchanging scalar values. The eigenvector is computed once before the bandit phase and stored locally.

\section{Max-Flooding and Routing Tree}
\label{app:flooding}

\subsection{Routing Targets}

Given the approximate eigenvector $\tilde{\psi}$, each agent $i$ selects a routing target:
\[
j_i^\star = \arg\max_{j \in \mathcal{N}(i)} \tilde{\psi}_j, \quad \text{subject to } \tilde{\psi}_{j_i^\star} > \tilde{\psi}_i.
\]
If no neighbor has strictly higher $\tilde{\psi}$ than $i$, then $i$ is a \emph{local maximum}. Each agent needs only its own $\tilde{\psi}_i$ and the values $\tilde{\psi}_j$ of its immediate neighbors to determine $j_i^\star$, a purely local computation.

\subsection{Max-Flooding to a Single Tree}

The local routing rule above may produce multiple local maxima when the graph has several structurally separated dense regions. To merge all local maxima into a single spanning tree rooted at the global maximum $h^\star = \arg\max_i \tilde{\psi}_i$, we run a short \emph{max-flooding} phase after power iteration, before the bandit phase begins.

\paragraph{Max-flooding protocol.}
Each agent $i$ maintains a flooded value $m_i^{(0)} = \tilde{\psi}_i$ and a pointer $\mathrm{via}_i$ initialized to $i$ itself. The protocol runs for a fixed number of rounds $T_{\mathrm{flood}}$, a hyperparameter set to any value at least $\mathrm{diam}(G)+1$:
\begin{enumerate}
    \item Each agent $i$ broadcasts $m_i^{(\tau-1)}$ to all neighbors $j \in \mathcal{N}(i)$.
    \item Agent $i$ updates: if $m_j^{(\tau-1)} > m_i^{(\tau-1)}$ for some neighbor $j$, set $m_i^{(\tau)} \leftarrow m_j^{(\tau-1)}$ and $\mathrm{via}_i \leftarrow j$; otherwise keep $m_i^{(\tau)} \leftarrow m_i^{(\tau-1)}$.
    \item After round $\tau = T_{\mathrm{flood}}$, every agent stops flooding and proceeds to the first bandit cycle.
\end{enumerate}
Because the global maximum propagates one hop per round, $\mathrm{diam}(G)+1$ rounds suffice for it to reach every node, so any $T_{\mathrm{flood}} \geq \mathrm{diam}(G)+1$ guarantees convergence. No agent needs to know $\mathrm{diam}(G)$ or any other global property of the graph; each agent needs only the shared round budget $T_{\mathrm{flood}}$, treated as an ordinary hyperparameter. A fixed common budget is what a purely local, silence-based stopping rule cannot provide: since agents at different distances from $h^\star$ stop updating at different rounds, a per-agent stopping rule would let them enter the first pull phase at different rounds, breaking the synchronized-cycle structure the relay protocol depends on. A common $T_{\mathrm{flood}}$ on the synchronous clock instead has every agent finish flooding on the same round and begin the first pull together.

After convergence, $m_i = \max_{j \in V} \tilde{\psi}_j$ for all $i$, and $\mathrm{via}_i$ points toward the neighbor through which the global maximum was first reached. Every local maximum $\ell$ checks locally whether $m_{\mathrm{via}_\ell} > \tilde{\psi}_\ell$; if so, it re-routes by setting $j_\ell^\star \leftarrow \mathrm{via}_\ell$. The unique node $h^\star$ whose own $\tilde{\psi}_{h^\star}$ equals the global maximum never sees a better value from any neighbor, so it keeps $j_{h^\star}^\star$ undefined and becomes the root. Since each re-routing step moves strictly toward the global maximum, no cycles can form.

The result is a single directed spanning tree $\mathcal{T}$ rooted at $h^\star$, with all edges oriented from children toward the root.

\begin{remark}
Max-flooding runs for $T_{\mathrm{flood}}$ rounds, each exchanging a single scalar per edge. The termination round is fixed in advance and shared, requiring no per-agent stopping decision and no global coordination beyond the common hyperparameter. On well-connected graphs the diameter is $O(\log N)$, so a budget of $T_{\mathrm{flood}} = O(\log N)$ suffices, making this a negligible overhead relative to the bandit horizon $T$.
\end{remark}

\subsection{Routing Tree Definitions}

The routing targets, after max-flooding, define a single directed spanning tree $\mathcal{T}$ on $V$ with unique root $h^\star$. Define:
\begin{itemize}
    \item $d_i$: the depth of node $i$ in $\mathcal{T}$, i.e., the number of hops from $i$ to $h^\star$ along the routing path.
    \item $D = \max_{i \in V} d_i$: the height of $\mathcal{T}$.
    \item $\mathrm{par}(i)$: the parent of $i$ in $\mathcal{T}$ (i.e., $j_i^\star$); undefined for $h^\star$.
    \item $\mathrm{ch}(i)$: the set of children of $i$ (nodes $j$ with $j_j^\star = i$).
\end{itemize}

\section{Relay Cycle Illustration}
\label{app:relay_cycle}

Figure~\ref{fig:relay_cycle} illustrates one full relay cycle on a depth-2 tree ($D=2$, cycle length $2D+1=5$ rounds). The hub is at the top; two intermediate nodes at depth 1 each have two leaf children at depth 2.

\begin{figure}[H]
    \centering
    \includegraphics[width=\textwidth]{results/MERW_visualization/relay_cycle.pdf}
    \caption{One relay cycle on a depth-2 tree ($N=7$, $D=2$). From left to right: $t=0$ -- all nodes pull their arm locally; $t=1$ -- every node sends its packet one hop toward the hub; $t=2$ -- intermediate nodes relay the packet to the hub; $t=3$ -- hub runs UCB and sends the updated snapshot one hop down; $t=4$ -- intermediate nodes forward the snapshot to their leaves, which schedule their next pull.}
    \label{fig:relay_cycle}
\end{figure}

\section{Full Algorithm Pseudocode}
\label{app:pseudocode}

Regret minimization and best arm identification are both instances of a single generic protocol, \merw{} (Algorithm~\ref{alg:merw-generic}): after the shared setup phase (distributed power iteration and max-flooding, Appendices~\ref{app:poweriter}--\ref{app:flooding}), every node initializes its local state and every cycle thereafter repeats the same four phases -- \textbf{Initialization}, \textbf{Pull}, \textbf{Uplink} ($D$ rounds), \textbf{Hub decision} (1 round, hub-local), \textbf{Downlink} ($D$ rounds) -- with three application-specific hooks left abstract: $\textsc{Init}$, $\textsc{LocalPull}$, and $\textsc{HubDecide}$. Table~\ref{tab:instantiations} instantiates all three hooks for the two applications studied in this paper; any other decentralized cooperative-learning task whose decision rule depends only on globally aggregated per-agent statistics (in the sense of Theorem~\ref{thm:coordinator}) can be run over the same tree by supplying a different set of hooks.

\begin{algorithm}[H]
\caption{\merw{} (generic protocol)}
\label{alg:merw-generic}
\begin{algorithmic}[1]
\Require Graph $G$; application-specific hooks $\textsc{Init}$, $\textsc{LocalPull}$, $\textsc{HubDecide}$; stopping condition on $t$ or on a hub-computed flag $\mathrm{done}$
\State Compute $\tilde{\psi}$ via distributed power iteration (Appendix~\ref{app:poweriter})
\State Run max-flooding and re-route local maxima toward the global maximum (Appendix~\ref{app:flooding})
\State Single spanning tree $\mathcal{T}$ with root $h^\star$; each node $i$ knows only its own $\mathrm{par}(i)$ so far -- $\mathrm{ch}(i)$, $d_i$, and $D$ are discovered over the first cycle below
\State $t \leftarrow 0$;\quad $\mathrm{done} \leftarrow \textbf{false}$ \Comment{$t$: wall-clock round, tracked locally by every node}

\State \textcolor{gray}{$\triangleright$ \textbf{Initialization Phase}}
\ForAll{$i \in V$}
    \State $\sigma_i \leftarrow \textsc{Init}()$;\quad $t \mathrel{+}= (\text{rounds consumed by this call})$
\EndFor

\Repeat
    \State \textcolor{gray}{$\triangleright$ \textbf{Pull phase}}
    \ForAll{$i \in V$}
        \State $(b_i, \sigma_i) \leftarrow \textsc{LocalPull}(\sigma_i)$
        \State $t \mathrel{+}= (\text{rounds consumed by this call})$
    \EndFor

    \State \textcolor{gray}{$\triangleright$ \textbf{Uplink}}
    \ForAll{$i \in V \setminus \{h^\star\}$} send $b_i$ to $\mathrm{par}(i)$;\quad $b_i \leftarrow \mathbf{0}$ \EndFor
    \Repeat
        \State $t \mathrel{+}= 1$;\quad hub absorbs any packets received this round into $B$
        \ForAll{$i \in V \setminus \{h^\star\}$}
            \If{$i$ received a packet from a child this round}
                \State $b_i \mathrel{+}= (\text{sum of packets received})$;\quad send $b_i$ to $\mathrm{par}(i)$;\quad $b_i \leftarrow \mathbf{0}$ 
            \EndIf
        \EndFor
    \Until{hub receives nothing this round} 

    \State \textcolor{gray}{$\triangleright$ \textbf{Hub decision} ($t \mathrel{+}= 1$, hub-local)}
    \State $(m, \mathrm{done}) \leftarrow \textsc{HubDecide}(B)$ \Comment{$m$: the message broadcast on the downlink (running aggregate, surviving set, etc.); $\mathrm{done}$ uses no information the hub doesn't already have}
    \State $t \mathrel{+}= 1$

    \State \textcolor{gray}{$\triangleright$ \textbf{Downlink} ($D$ rounds -- $D$ is known to the hub, having just observed the uplink go silent for one round; every other node learns $D$ and its own depth $d_i$ from this broadcast, not before}
    \State hub sends $(m,\, D,\, \mathrm{done},\, \mathrm{hops}{=}0)$ to each child in $\mathrm{ch}(h^\star)$ \Comment{the hub now also knows $\mathrm{ch}(h^\star)$, having just heard from them in the Uplink}
    \Repeat
        \State $t \mathrel{+}= 1$
        \ForAll{$i \in V \setminus \{h^\star\}$}
            \If{$i$ received $(m, D, \mathrm{done}, \mathrm{hops})$ this round}
                \State $d_i \leftarrow \mathrm{hops} + 1$;\quad update $\sigma_i$ from $m$;\quad forward $(m, D, \mathrm{done}, \mathrm{hops}{+}1)$ to each child in $\mathrm{ch}(i)$ \Comment{purely local: $i$ only ever reacts to the packet it just received, and $\mathrm{ch}(i)$ is whichever nodes route through $i$}
            \EndIf
        \EndFor
    \Until{$D$ rounds have elapsed since the hub sent} \Comment{$i$ never independently evaluates the stopping condition; it only learns $\mathrm{done}$ from the packet}
\Until{$\mathrm{done}$ \textbf{or} $t \geq T$ (whichever stopping condition applies)}
\end{algorithmic}
\end{algorithm}

\begin{table}[H]
\centering
\caption{Instantiating Algorithm~\ref{alg:merw-generic} for the two applications of this paper.}
\label{tab:instantiations}
\begin{tabular}{p{0.19\textwidth}p{0.36\textwidth}p{0.36\textwidth}}
\toprule
 & \merw-UCB (Section~\ref{sec:regret}) & \merw-BAI (Section~\ref{sec:bai}) \\
\midrule
$\textsc{Init}$ & pull every arm $k \in \{1,\ldots,K\}$ once; $\sigma_i \leftarrow (\hat n_i, \hat s_i)$, $K$ rounds & $\sigma_i \leftarrow S_0 = \{1,\ldots,K\}$, 0 rounds (no pulls needed) \\
$\textsc{LocalPull}$ & pull $a_i = \arg\max_k[\hat s_i^k/\hat n_i^k + \sqrt{c\log \max(P,1)/\hat n_i^k}]$; 1 round & pull every $k \in S_{r-1}$ exactly $L_r$ times; $L_r$ rounds ($L_r$, $\epsilon_r$ pre-agreed from $N,K,\delta$) \\
Uplink payload $b_i$ & new $(\hat n_i, \hat s_i)$ observations & local sums $\{s_k^{(i)}\}_{k \in S_{r-1}}$ \\
$\textsc{HubDecide}$ & $P \mathrel{+}= N$; pull $a_h$ via UCB on $B=(\hat n,\hat s)$; $P\mathrel{+}=1$; never stops early ($\mathrm{done} = (t \geq T)$) & $\tilde p_k \leftarrow B_k/(N L_r)$; $S_r \leftarrow \{k : \tilde p_k \geq \max_j \tilde p_j - \epsilon_r\}$; $\mathrm{done} = (|S_r|=1)$; hub never pulls \\
Downlink payload $m$ & full aggregate $(\hat n, \hat s, P)$ & surviving set $S_r$ only \\
Cycle length & fixed, $2D+1$ & varies, $L_r + 2D$ (shrinks as $\epsilon_r$ tightens) \\
\bottomrule
\end{tabular}
\end{table}

The two instantiations differ only in these five rows; the surrounding relay mechanics -- the hop-by-hop uplink fold, the hub-only stopping check, and the hop-by-hop downlink broadcast -- are identical, since they are exactly the aggregation and broadcast primitives of Propositions~\ref{prop:aggregation}--\ref{prop:broadcast} that make Theorem~\ref{thm:coordinator} apply to any hook pair whose $\textsc{HubDecide}$ depends only on the aggregated $B$.

\section{Additional Regret Results}
\label{app:regret}

Figures~\ref{fig:regret_barbell}--\ref{fig:regret_grid} show cumulative regret on the remaining graph families. The three-panel layout is identical to Figure~\ref{fig:regret_ba}: group regret vs.\ round, hub regret vs.\ round, and group regret vs.\ total arm pulls.

On the barbell graph the hub concentrates in one of the two cliques; nodes in the other clique route through the bridge edge, resulting in a larger $D$ and correspondingly slower per-round progress visible in the flat segments of the wall-clock regret curve. On the lollipop, the hub sits at the clique core and the path nodes form a deep chain ($D$ up to $N/2$), making the flat segments especially pronounced. On the chain the hub is the central node and $D = N/2$, so each cycle is long; nevertheless regret vs.\ total pulls remains well-behaved. The star is the ideal case: $D=1$, every cycle lasts 3 rounds, and the hub aggregates all observations with minimal delay. On the grid MERW selects the center node and $D$ equals the grid radius, giving moderate cycle length.

Across all graph families the regret-vs.-pulls curves are consistent, confirming that the sample efficiency of \merw{} is independent of graph topology: the relay structure adds communication overhead (visible in wall-clock regret) but does not waste observations.

\begin{figure}[H]
    \centering
    \includegraphics[width=\textwidth]{results/Regret/relay_regret_barbell_N20_K5_T5000.pdf}
    \caption{Barbell graph ($N=20$, $K=5$, $T=5000$, 50 runs).}
    \label{fig:regret_barbell}
\end{figure}

\begin{figure}[H]
    \centering
    \includegraphics[width=\textwidth]{results/Regret/relay_regret_lollipop_N20_K5_T5000.pdf}
    \caption{Lollipop graph ($N=20$, $K=5$, $T=5000$, 50 runs).}
    \label{fig:regret_lollipop}
\end{figure}

\begin{figure}[H]
    \centering
    \includegraphics[width=\textwidth]{results/Regret/relay_regret_chain_N20_K5_T5000.pdf}
    \caption{Chain graph ($N=20$, $K=5$, $T=5000$, 50 runs).}
    \label{fig:regret_chain}
\end{figure}

\begin{figure}[H]
    \centering
    \includegraphics[width=\textwidth]{results/Regret/relay_regret_star_N20_K5_T5000.pdf}
    \caption{Star graph ($N=20$, $K=5$, $T=5000$, 50 runs).}
    \label{fig:regret_star}
\end{figure}

\begin{figure}[H]
    \centering
    \includegraphics[width=\textwidth]{results/Regret/relay_regret_grid_N25_K5_T5000.pdf}
    \caption{Grid graph ($N=25$, $K=5$, $T=5000$, 50 runs).}
    \label{fig:regret_grid}
\end{figure}

\end{document}
